\section*{Ex.~2.8 (22.5 pts) --- Three Papers on Scheduling, Found by Search}

\emph{How I found them.} I searched for highly cited scheduling papers outside the
course paper list whose full text is freely downloadable, then picked one from each of
three eras --- 1973 uniprocessor real-time theory, 2016 multicore kernel, 2019
datacentre microsecond-scale --- so the three critiques would bite on different
evidence rather than repeat one argument.

\subsection*{Paper 1 --- Liu and Layland, ``Scheduling Algorithms for
Multiprogramming in a Hard-Real-Time Environment'' (JACM, 1973)}

\paragraph{Summary.}
The paper asks how much of one processor can be committed to $m$ periodic tasks while
still guaranteeing every deadline. A task is reduced to two numbers, its period $T_i$
and its run-time $C_i$, under five assumptions: periodic requests, a deadline equal to
the next request, independent tasks, constant run-times, and aperiodic work confined to
initialisation and failure recovery.

For fixed priorities, Theorem~1 identifies the \emph{critical instant} --- ``A critical
instant for any task occurs whenever the task is requested simultaneously with requests
for all higher priority tasks'' --- which turns an infinite schedule into a finite
check. Theorem~2 shows the \emph{rate-monotonic} rule (shorter period, higher priority,
regardless of run-time) is optimal among all fixed-priority rules, and Theorem~5 bounds
utilisation at $U = m\left(2^{1/m}-1\right)$: about $0.83$ for two tasks, $0.78$ for
three, falling to $\ln 2 \approx 0.69$. Section~7 instead assigns priorities by nearest
deadline, and Theorem~7 proves that \emph{deadline driven} rule feasible if and only if
$\sum_i C_i/T_i \le 1$, so the same hardware is schedulable at 100\,\% load. Section~8
mixes the two, running the $k$ shortest-period tasks rate-monotonically on the
interrupt hardware; the worked example in Section~9, $T = (3,4,5)$ with $C_1 = C_2 = 1$,
reaches $98.3\,\%$ mixed against $100\,\%$ deadline-driven and $78.3\,\%$
rate-monotonic.

\paragraph{Critical judgment.}
The cost model is the weak point. Assumption (A4) folds context-switch cost into the
constant $C_i$, claiming ``the bookkeeping time necessary to request a successor and the
costs of preemptions can be taken into account''. It cannot: how often a task is
preempted depends on the whole task set, not on the task, and rate-monotonic scheduling
preempts far more than deadline-driven scheduling does. The practical case for the
fixed-priority algorithm is exactly the quantity the model prices at zero. Section~6
concedes that ``the practical costs of switching between tasks must still be counted'',
and then never counts them.

The more consequential gap is independence. The conclusion is candid about the other
assumptions, calling periodicity and constant run-times ``perhaps the most important and
least defensible'' and admitting that without periodicity ``None of our analytic work
would remain valid''; it even reopens the deadline assumption. Independence is examined
once, and only in its easy case: (A3) ``does not exclude the situation in which the
occurrence of a task $\tau_2$ can only follow a certain (fixed) number\ldots of
occurrences of a task $\tau_1$'', which is absorbed by choosing commensurate periods.
Fixed request-count coupling is not what real systems break on. They share mutexes, and a
low-priority task holding a lock blocks a high-priority one for a time the utilisation
bound has no term for. That hole stayed open until Sha, Rajkumar and Lehoczky's
priority-inheritance protocols added an explicit blocking term to the test; the 1997
Mars Pathfinder resets showed what it costs in the field. The framework itself won
completely --- but the paper's parting advice, to make periodicity and constant
run-times ``a design goal'', is the part practice discarded, in favour of sporadic-task
and server models.

\subsection*{Paper 2 --- Lozi, Lepers, Funston, Gaud, Qu\'ema and Fedorova,
``The Linux Scheduler: a Decade of Wasted Cores'' (EuroSys, 2016)}

\paragraph{Summary.}
The paper takes one invariant --- ``make sure that ready threads are scheduled on
available cores'' --- and shows Linux's Completely Fair Scheduler violating it for
seconds at a time. The cause is structural: per-core runqueues force periodic load
balancing, balancing is too expensive to run often, and the load metric (weight scaled
by utilisation, then divided by the cgroup's thread count) is complicated enough to hide
imbalance.

Four bugs follow. \emph{Group Imbalance}: a stealing core compares the \emph{average}
load of a scheduling group, so one heavy thread masks idle cores beside it; comparing
minimum load instead made a kernel \texttt{make} 13\,\% faster and NAS \texttt{lu}
$13\times$ faster. \emph{Scheduling Group Construction}: groups are built from core~0's
viewpoint, so nodes two hops apart land in both of the groups $\{0,1,2,4,6\}$ and
$\{1,2,3,4,5,7\}$, their averages always match and no work moves --- \texttt{lu} takes
1040\,s instead of 38\,s. \emph{Overload-on-Wakeup}: for cache reuse a waking thread
is placed only among cores of the waker's node, so it can wake onto a busy core while
others idle; the fix cut TPC-H query~18 by $22.2\,\%$ and the full benchmark by
$13.2\,\%$. \emph{Missing Scheduling Domains}: a call dropped during refactoring means
disabling and re-enabling any core silently stops all cross-NUMA balancing, and
\texttt{lu} runs $137.59\times$ slower. None of these crash anything, so the authors add
an online sanity checker for the invariant and a trace visualiser, the latter needing
under 150 lines of kernel instrumentation.

\paragraph{Critical judgment.}
The bug analyses are exact, but the evidence is wide in workloads and narrow in
hardware: every measurement comes from one 8-node AMD Bulldozer server. Table~4
nevertheless records ``Impacted applications: All'' for three of the four bugs, two of
them topology-dependent --- Scheduling Group Construction needs nodes two hops apart ---
so a claim about which \emph{applications} suffer is doing duty as one about which
\emph{machines} do, on a sample of one.

The energy claim is weaker. The introduction asserts that ``Energy waste is
proportional'' to the lost performance, yet power runs through the paper as a design
pressure without ever being measured --- no power figure, axis or table --- and it bites
on the Overload-on-Wakeup fix. The authors raise that tradeoff themselves rather than
hiding it, since waking a thread on a long-idle core forces it out of a low-power state,
but
their answer is to ``only enforce the new wakeup strategy if the system's power
management policy does not allow cores to enter low-power states at all''. The fix that
wins most on their database workload is therefore disabled on virtually every production
server, and the energy argument that might justify re-enabling it is asserted rather
than measured.

History has been kind to both forecasts. ``New scheduler designs come and go'' held:
CFS was itself replaced, by EEVDF, in kernel 6.6 (2023). The modular alternative of
Section~5 --- ``a scheduler that is a collection of modules: the core module and
optimization modules'' --- was then taken up in spirit by \texttt{sched\_ext} in 6.12
(2024), which moves policy into BPF programs outside the core. The fit is loose, since
\texttt{sched\_ext} replaces the whole scheduler rather than advising an
invariant-keeping core, but the instinct is the paper's.

\subsection*{Paper 3 --- Kaffes, Chong, Humphries, Belay, Mazi\`eres and Kozyrakis,
``Shinjuku: Preemptive Scheduling for $\mu$second-scale Tail Latency'' (NSDI, 2019)}

\paragraph{Summary.}
Dataplane operating systems reach microsecond latencies by never preempting: IX hashes
requests to per-core queues (distributed FCFS), and ZygOS work-steals to approximate a
central FCFS queue. Both let a short request queue behind a long one, which is fatal
when service times are heavy-tailed or bimodal, where theory prefers processor sharing
--- and processor sharing needs preemption. Linux cannot supply it, since CFS targets a
6\,ms preemption latency with a 0.75\,ms minimum.

Shinjuku is a single-address-space OS on Dune in which one dispatcher thread owns the
queues and interrupts workers directly. Removing both VM exits from the interrupt path
--- posted
interrupts on the receiver, a directly mapped APIC on the sender --- takes the sender
from 2081 to 298 cycles and the receiver from 2662 to 1212, and context switches drop to
36--109 cycles; it is that 1212-cycle receiver cost that makes preemption every
5\,$\mu$s affordable, at under $10\,\%$ of worker throughput. Two policies exploit it: a
single queue preempting every 5--15\,$\mu$s, and a multi-queue policy serving whichever
queue's head request has the largest ratio of queueing time to its SLO target. Against
both IX and ZygOS, Shinjuku obtains up to $50\,\%$ lower tail latency \emph{at low load}
and $5\times$ the throughput for a given 300\,$\mu$s tail-latency target on a bimodal
synthetic workload, and on RocksDB serving 99.5\,\% \texttt{GET} plus 0.5\,\%
\texttt{SCAN}, $88\,\%$ lower tail latency with $6.6\times$ the throughput.

\paragraph{Critical judgment.}
Centralisation is what makes the policies possible and is the claim least tested.
Section~3.6 admits that one dispatcher ``comfortably saturate[s] a full socket'' and no
more --- \S4.3 measures it scaling near-linearly to 11 workers at 5\,M requests/s, and
9.5\,M with two dispatchers --- and proposes to scale further by ``steer[ing] requests to
different dispatchers using the NIC RSS feature''. That is the
distributed queueing Section~2 blames for head-of-line blocking, reintroduced at the
granularity of dispatcher groups, and what it costs is never measured: the only
multi-dispatcher experiment issues ``short requests with 1$\mu$sec fixed service time to
stress the dispatcher'', the one distribution under which head-of-line blocking cannot
occur. The headline claim is about heavy tails; the scaling result is measured where
tails are absent.

The security discussion is honest, but its remedy is still a promise. Mapping APIC
registers into the process means, as the authors say, that ``one process could launch a
denial-of-service attack on another process by issuing a large number of interrupts to a
specific core''; they answer that ``future versions of Shinjuku will run application
code in ring 3 $\ldots$ eliminating this attack vector with very small overhead'', citing
an 84-cycle ring transition. But every evaluation number was taken with applications in
ring~0, chosen ``to avoid the address space crossings'' --- crossings that, at a
5\,$\mu$s quantum, would recur on every preemption. The system that was measured is not
the system that would be deployable.

The premise --- that the obstacle was mechanism rather than policy --- has held up: the
hardware the paper asks for, ``support for lightweight user-level interrupts'', later
shipped as Intel's user interrupts (UINTR) in Sapphire Rapids.

\subsection*{References}

\begin{description}
\item[{[1]}] C.~L. Liu and J.~W. Layland. ``Scheduling Algorithms for Multiprogramming
  in a Hard-Real-Time Environment.'' \emph{Journal of the ACM}, vol.~20, no.~1,
  January 1973, pp.~46--61. DOI: 10.1145/321738.321743.
\item[{[2]}] J.-P. Lozi, B.~Lepers, J.~Funston, F.~Gaud, V.~Qu\'ema and A.~Fedorova.
  ``The Linux Scheduler: a Decade of Wasted Cores.'' \emph{Proc. 11th European
  Conference on Computer Systems (EuroSys~'16)}, London, April 2016.
  DOI: 10.1145/2901318.2901326.
\item[{[3]}] K.~Kaffes, T.~Chong, J.~T. Humphries, A.~Belay, D.~Mazi\`eres and
  C.~Kozyrakis. ``Shinjuku: Preemptive Scheduling for $\mu$second-scale Tail Latency.''
  \emph{Proc. 16th USENIX Symposium on Networked Systems Design and Implementation
  (NSDI~'19)}, Boston, February 2019, pp.~345--360.
  \texttt{https://www.usenix.org/\allowbreak conference/\allowbreak
  nsdi19/\allowbreak presentation/\allowbreak kaffes}
\item[{[4]}] L.~Sha, R.~Rajkumar and J.~P. Lehoczky. ``Priority Inheritance Protocols:
  An Approach to Real-Time Synchronization.'' \emph{IEEE Transactions on Computers},
  vol.~39, no.~9, September 1990, pp.~1175--1185. (Cited above for the blocking term
  missing from~[1].)
\end{description}
